%% kapa: exam=1 round=30 year=2019 period=1 form=A subject=civil_law subject_ko=민법 seq=2 number=2
%% source: Q-Net 감정평가사 기출문제 (cst00309) artlSeq=5205205 file=extracted/1008/5205205/2019년 제30회 감정평가사 1차시험 문제지/2019년 제30회 감정평가사 1차시험 문제지/2019년도 제30회 감정평가사 1차 1교시 A형.hwp
%% conversion: kordoc 4.24.51 -> tools/gen1.py (auto); review=proofread-sonnet; flags=none
\begin{question}{2}
민법상 권리에 관한 설명으로 옳지 않은 것은?
\begin{choices}
  \item 조건부권리는 기대권에 속한다.
  \item 채권과 청구권은 동일한 개념이다.
  \item 지상권자의 지료증감청구권은 형성권이다.
  \item 보증인의 최고ㆍ검색의 항변권은 연기적 항변권이다.
  \item 주된 권리가 시효로 소멸하면 종된 권리도 소멸한다.
\end{choices}
\end{question}
\answer{2}
